\subsection*{Dense radicals, ordinary resolvent collapse, and a bounded-error reduction}

All operators in this subsection are the complete physical Weil operators,
not the corrected logarithmic generators in a Weil-metric completion.
Put $H=L^2(\mathbb R,dx)$, $I_a=(-a,a)$, $\lambda=e^a$, and let
$J_a:L^2(I_a)\to H$ be zero extension. Inner products remain linear in
the first slot. Write $q_a$ for the canonical closed form of
$\mathsf W_{e^a}$. On compactly supported smooth tests it is the
restriction of a single support-consistent geometric form $q$.
The form domain and both pole signs are those of
Propositions~\ref{prop:v118-log-dirichlet} and
\ref{prop:v135-both-parities}. No Fourier cutoff is imposed.

\begin{lemma}[A total family of localized ordinary near-radicals]
\label{lem:v136-total-radicals}
Use the normalized source $\phi$ and smooth physical cutoff $\chi_a$
from Proposition~\ref{prop:v135-growing-radical}. Then
\[
 \mathcal R=\operatorname{span}\{\phi(\cdot-t):t\in\mathbb R\}
 \quad\hbox{is dense in }H.
\]
For every fixed $h\in\mathcal R$, $g_a=\chi_a h$ belongs to
$C_c^\infty(I_a)$ and satisfies
\begin{equation}\label{eq:v136-dense-residual}
 J_ag_a\longrightarrow h,\qquad
 \norm{\mathsf W_{e^a}g_a}_2\longrightarrow0.
\end{equation}
The even and odd reflected translate spans are dense in their respective
parity subspaces. These are ordinary-norm assertions; they assert neither
uniform approximation coefficients nor density in a global form norm.
\end{lemma}
\begin{proof}
The double-exponential estimate \eqref{eq:v135-gaussian-tail} makes
$\widehat\phi$ entire. It is not identically zero, so its real zeros
have measure zero. If $v\in H$ is orthogonal to every translate, the
inverse Fourier transform of the $L^1$ function
$\widehat v\,\overline{\widehat\phi}$ vanishes. Fourier uniqueness
gives $\widehat v\,\overline{\widehat\phi}=0$ almost everywhere,
and hence $v=0$. This is the elementary Hilbert-space version of the
classical translation-density principle~\cite{Wiener1932}.
Apply the bounded parity projections to obtain the two parity statements.

Write $h=\sum_{j=1}^r c_j\phi(\cdot-t_j)$ with fixed real $t_j$.
For sufficiently large $a$, every $|t_j|\le a/2$. The exterior-tail
proof of Proposition~\ref{prop:v135-growing-radical} uses only that
inequality, not membership in the spaced lattice. It therefore gives
\[
 \norm{\mathsf W_{e^a}(\chi_a\phi(\cdot-t_j))}_2
 \le C\sqrt a\,e^{-e^a/100}.
\]
The triangle inequality proves the residual assertion with a constant
depending on this fixed representation of $h$. Dominated convergence
proves the ordinary limit. In particular all global prime rows in the
exterior remainder, both poles and the logarithmic diagonal are retained.
\end{proof}

\begin{proposition}[Ordinary strong-resolvent collapse]
\label{prop:v136-resolvent-collapse}
Define the self-adjoint operator
\[
 A_a=\mathsf W_{e^a}\oplus0
 \quad\hbox{on }H=L^2(I_a)\oplus L^2(I_a^c),
 \qquad
 \mathcal D(A_a)=\mathcal D(\mathsf W_{e^a})\oplus L^2(I_a^c).
\]
Then, without a uniform lower bound or a sign hypothesis,
\begin{equation}\label{eq:v136-resolvent-collapse}
 (A_a-z)^{-1}\longrightarrow-z^{-1}I
 \quad\hbox{strongly for every }z\in\mathbb C\setminus\mathbb R.
\end{equation}
If $P_a$ is an orthogonal projection with range in $\mathcal D(A_a)$
and $\norm{A_aP_a}\to0$, the same conclusion holds for the operator
which deletes its block and assigns zero on its range. Precisely, define
\[
 A_a^{\perp}=A_a-A_aP_a-P_aA_a+P_aA_aP_a
 \quad\hbox{on }\mathcal D(A_a).
\]
The terms involving $P_aA_a$ mean their bounded extensions.
This applies to the actual repaired rank-one source and to the growing
block of Proposition~\ref{prop:v135-growing-radical}.
\end{proposition}
\begin{proof}
For $h\in\mathcal R$ use $g_a$ from Lemma~\ref{lem:v136-total-radicals},
regarded as a vector in $H$. The resolvent bound for self-adjoint
operators and the resolvent identity give
\begin{align*}
 \norm{(A_a-z)^{-1}h+z^{-1}h}
 &\le\left(\frac1{|\Im z|}+\frac1{|z|}\right)\norm{h-g_a}
       +\frac{\norm{A_ag_a}}{|z|\,|\Im z|}.
\end{align*}
This tends to zero. Uniform resolvent boundedness and density extend
the conclusion to every $h\in H$. For the deleted block, the difference
is a bounded self-adjoint perturbation of norm at most
$3\norm{A_aP_a}$. The resolvent identity bounds the new resolvent
difference by $3\norm{A_aP_a}/|\Im z|^2$.
\end{proof}

This concerns the ordinary $L^2$ realization of the Weil operator.
It does not identify a limit of the logarithmic generators in the
different Weil-metric spaces discussed, for example, in
\cite[Section 7]{Suzuki2026Screw}. Nor does it give convergence of
$q_a[f]$ for fixed compactly supported $f$; those values are exactly
support-consistent already. Strong-resolvent convergence alone does
not control unbounded spectral moments or lower spectral edges.
Indeed, on $\ell^2$, the operators
$T_n=-n|e_n\rangle\langle e_n|$ converge to zero in this topology,
while $\inf\sigma(T_n)=-n$. Even replacing $n$ by $1$ leaves a
fixed negative edge. These are operator-topology countermodels,
not arithmetic Weil tests.

Even support consistency and a dense radical do not by themselves
fix the sign. On $c_{00}\subset\ell^2$ take
$q_\pm[x]=\pm|\sum_kx_k|^2$; its $n$-coordinate matrices are
$\pm\mathbf1_n\mathbf1_n^*$. Their common limiting radical,
the finite sequences of sum zero, is dense: for example
$e_1-m^{-1}\sum_{k=2}^{m+1}e_k\to e_1$.
Both zero-extended matrix families converge strong-resolvent to zero,
since their unit nonzero-eigenvalue direction
$n^{-1/2}\mathbf1_n$ tends weakly to zero. The positive family is
nonnegative and the negative family has lower edge $-n$.
The latter amplifies the form value $-1$ to Rayleigh value $-m$
on $m^{-1}\sum_{k=2}^{m+1}e_k$.

\begin{proposition}[A uniform finite lower floor suffices]
\label{prop:v136-bounded-floor}
Suppose that for a cofinal sequence $a_j\to\infty$ there is one
finite constant $C_*$, independent of $j$, such that
\begin{equation}\label{eq:v136-bounded-floor}
 q_{a_j}[f]\ge-C_*\norm f_2^2
 \quad\hbox{for all }f\in\mathcal D(q_{a_j}).
\end{equation}
Then $q[f]\ge0$ on every compact smooth test. Consequently every
complete finite-window form is nonnegative, and the Weil criterion
implies RH. Equivalently, if any compact smooth test has negative
Weil form, then
\begin{equation}\label{eq:v136-negative-divergence}
 \inf\sigma(\mathsf W_{e^a})\longrightarrow-\infty
 \qquad(a\longrightarrow\infty).
\end{equation}
This is a conditional reduction: no bound \eqref{eq:v136-bounded-floor}
is established here.
\end{proposition}
\begin{proof}
Suppose a compact smooth $f$ has $q[f]=-d<0$. Given $\varepsilon>0$,
choose a fixed $h\in\mathcal R$ with $\norm{f-h}<\varepsilon/2$.
Put $g_a=\chi_a h$ and $w_a=f-g_a$ for sufficiently large windows.
By support consistency and self-adjointness,
\[
 |q_a[w_a]-q[f]|
 \le (2\norm f+\norm{g_a})\norm{\mathsf W_{e^a}g_a}
 \longrightarrow0,\qquad
 \norm{w_a}\longrightarrow\norm{f-h}<\varepsilon/2.
\]
Thus eventually $q_a[w_a]<-d/2$ and $0<\norm{w_a}<\varepsilon$,
so the variational principle gives
\begin{equation}\label{eq:v136-negative-amplification}
 \inf\sigma(\mathsf W_{e^a})
 \le\frac{q_a[w_a]}{\norm{w_a}^2}
 <-\frac{d}{2\varepsilon^2}.
\end{equation}
Since $\varepsilon$ is arbitrary, this proves
\eqref{eq:v136-negative-divergence}, and contradicts every cofinal
finite lower floor. Therefore \eqref{eq:v136-bounded-floor} forces
$q[f]\ge0$ for every compact smooth $f$.
Compact smooth tests are a form core at each fixed window by
Proposition~\ref{prop:v118-log-dirichlet}, so complete local
positivity follows. The fixed-support Weil implication is the same
one used in Proposition~\ref{prop:v121-cofinal-rh}.
The argument gives no effective window threshold for a prescribed
$\varepsilon$: the finite translate coefficients and shifts can grow
as its ordinary approximation is refined.
\end{proof}

\begin{corollary}[Bounded-error version of the gap-free target]
\label{cor:v136-complement-floor}
Let $P_a$ be either the actual repaired source projection or the
growing radical projection, with its established
$\norm{\mathsf W_{e^a}P_a}\to0$. It is sufficient to prove, on a
cofinal family, the single uniform estimate
\begin{equation}\label{eq:v136-complement-target}
 q_a[f]\ge-C_*\norm f^2,
 \qquad f\in\mathcal D(q_a)\cap\operatorname{Ran}(I-P_a),
 \qquad C_*<\infty\text{ independent of }a.
\end{equation}
Both parity sectors must be included. In the signed concentration
criterion \eqref{eq:v131-weighted-concentration}, bounded cofinal
errors in both sectors therefore suffice; proving their decay to zero
is a stronger sufficient objective. This corollary does not identify
$C_a^{\rm low}$ with either displayed complement.
\end{corollary}
\begin{proof}
Proposition~\ref{prop:v118-gap-free} gives a full lower floor
$-C_*-(2/\sqrt3)\norm{\mathsf W_{e^a}P_a}$, which is uniformly
bounded along the family. Apply Proposition~\ref{prop:v136-bounded-floor}.
The even and odd blocks combine as in \eqref{eq:v135-full-error}.
For the physical deep/plunge decomposition, any missing block and
coupling still needs its own justified bound.
\end{proof}

\begin{proposition}[No persistent bounded positive comparison]
\label{prop:v136-positive-comparison}
Let $u_a$ be the actual ordinary-unit source, with $J_au_a\to\phi$
and $\norm{\mathsf W_{e^a}u_a}\to0$, and let
$Q_a=I-|u_a\rangle\langle u_a|$. Suppose
\[
 B_a=Q_aB_aQ_a\succeq0,\qquad\sup_a\norm{B_a}<\infty,
 \qquad
 q_a[f]\ge\ip{B_af}{f}-\eta_a\norm f^2\quad(f\perp u_a),
 \qquad\eta_a\to0,
\]
on the complete form domain. Then
$J_aB_aJ_a^*\to0$ strongly on $H$.
In particular a fixed nonzero positive bounded comparison on
$\phi^\perp$ cannot satisfy these assumptions.
If instead the comparison is imposed on the actual $C_a^{\rm low}$
and all its removed physical image columns are even, any uniformly
bounded positive comparison must tend strongly to zero on every
fixed odd vector. No identification of the two complements is used.
\end{proposition}
\begin{proof}
For $h\in\mathcal R$ set $f_a=Q_a(\chi_a h)$. The source assumptions
and Lemma~\ref{lem:v136-total-radicals} give
\[
 J_af_a\longrightarrow(I-|\phi\rangle\langle\phi|)h,
 \qquad\norm{\mathsf W_{e^a}f_a}\longrightarrow0.
\]
Consequently $\ip{B_af_a}{f_a}\to0$. If $M$ bounds $\norm{B_a}$,
positivity gives $\norm{B_af_a}^2\le M\ip{B_af_a}{f_a}\to0$.
The limiting projected translates are dense in $\phi^\perp$;
uniform boundedness extends the assertion to that subspace.
Compression and $J_au_a\to\phi$ give convergence on $\phi$ too.
For $C_a^{\rm low}$ use only odd reflected translates. Their odd
cutoffs belong to that physical complement exactly, and their span
is dense in the odd subspace. The same positive-square argument works
even when the comparison operator itself does not preserve parity.
\end{proof}

Uniform boundedness and the chosen topology matter: escaping rank-one
projections can have norm one and strong limit zero, while
$n^2|e_n\rangle\langle e_n|$ shows why testing a dense family without
a common norm bound is insufficient. A fixed closed nonnegative form
has a related obstruction only when all the projected cutoff tests
belong to its domain: lower semicontinuity then makes it vanish on
$\phi^\perp$, not necessarily on the source direction. This assertion
uses ordinary $L^2$ closedness and does not concern a Weil-metric
quotient completion.

The essential new project ingredient is totality of \emph{all fixed real
translates}, rather than a single uniformly separated growing lattice.
The classical translation-density and radical identities are not claimed
as new. These results strengthen the reduction and exclude persistent
positive comparisons; they do not supply an independent arithmetic
lower bound. The live signed concentration construction remains open,
now with the sufficient target \eqref{eq:v136-complement-target}.
No G2 sign gap is closed. The finite sampler, its explicit graph defect,
the physical endpoint and the existing Fourier scale factors are unchanged.
